\chapter{Returns of Financial Assets and Instruments}
\label{chap:Returns of Financial Assets and Instruments}

\begin{center}
\begin{tabular}{p{0.06\linewidth} | p{0.88\linewidth}}
	\toprule
	\multicolumn{2}{l}{\textbf{LEARNING OUTCOMES}} \\
	\midrule
	a & describe, compare, and interpret returns \\
	\midrule
	b & describe, compare, and interpret required rates of return, risk-free rates, risk premia, and inflation \\
	\bottomrule
\end{tabular}
\end{center}

\vspace{1em}

\begin{itemize}
	\item[] \textbf{Financial Returns}
	      \begin{itemize}
		      \item[] Total Return
		      \item[] Arithmetic and Geometric Holding Period Return
		      \item[] Compounding: Daily, Monthly, Quarterly, and Annual Returns
		      \item[] Continuous Compounding
	      \end{itemize}
	\item[] \textbf{Common Return Measures}
	      \begin{itemize}
		      \item[] The Required Rate of Return
		      \item[] Excess Returns
		      \item[] Real Returns
		      \item[] Gross and Net Return
		      \item[] Pre-Tax and After-Tax Nominal Return
		      \item[] Leveraged Returns
	      \end{itemize}
\end{itemize}

\section{FINANCIAL RETURNS}
\begin{center}
\begin{tabular}{p{0.06\linewidth} | p{0.88\linewidth}}
	a & describe, compare, and interpret returns \\
\end{tabular}
\end{center}

\textcolor{blue}{
	\textit{Lãi suất hay lợi suất là đối tượng thống kê chính trong hoạt động đầu tư. 
Có thể hiểu, thị trường tài chính là nơi buôn bán một loại hàng hóa đặc biệt, là Vốn.
Và lãi suất chính là mức giá cho loại hàng hóa đặc biệt đó.}}

\subsection{Total Return} % (fold)
\begin{equation}
\begin{alignedat}{2}
\text{Total Return}
    &={} \text{Capital Appreciation Return} &+{} & \text{Capital Distribution Return} \\
r   &={} r_{\text{price}} &+{} & r_{\text{distribution}} \\
    &={} \frac{P_1-P_0}{P_0} &+{} & \frac{Inc}{P_0}
\end{alignedat}
\end{equation}

For capital distributions on equities, we focus on the dividend yield; for bonds, we look at current yield.
\begin{itemize}
	\item Some expressions of Return:
	\begin{itemize}
		\item[] \textbf{Expected returns}, or ex ante returns, are the anticipated returns on an investment over a specified period.
		\item[] \textbf{Actual returns}, or ex post returns, are the returns on an investment after the investment period has ended.
		\item[] \textbf{Unrealized returns}: changes in value represent while an investor holds an investment.
		\item[] \textbf{Realized returns}: when the investment is ultimately sold, gains and losses would be realized at the current market value.
	\end{itemize}
\end{itemize}
% subsection Total Returndowncase:_}${1:_} (end)

\subsection{Arithmetic and Geometric Holding Period Return} % (fold)
% \label{sub:Arithmetic and Geometric Holding Period Return}
To ensure accurate comparisons across different investments and time periods, it's crucial that these returns are calculated using a uniform time period.
\begin{itemize}
	\item \textbf{arithmetic mean return}: 
	\begin{equation}
		\bar{r}_A = r_1 + r_2 + ... + r_{T-1} + r_T = \frac{1}{T}\sum_{t=1}^{T}r_t
		% \label{eq:}
	\end{equation}
	\item \textbf{geometric mean return}:
	\begin{equation}
		\bar{r}_G = \sqrt[T]{(1 + r_1) \times (1 + r_2) \times ... \times (1 + r_T)} - 1 = \sqrt[T]{\prod_{t = 1}^{T}(1 + r_t)} - 1
		% \label{eq:}
	\end{equation}
\end{itemize}
Geometric mean is a more accurate reflection of investment performance over time.
Arithmetic average return calculations assume that the amount invested at the beginning of the period does not change,
but the base amount changes due to capital appreciation or depreciation.
The arithmetic-geometric mean inequality highlights the impact of volatility on investment returns.
% subsection Arithmetic and Geometric Holding Period Return (end)

\subsection{Compounding: Daily, Monthly, Quarterly, and Annual Returns} % (fold)
% \label{sub:Compounding: Daily, Monthly, Quarterly, and Annual Returns}
With c is the number of periods in a year,
\begin{equation}
	r_{annualized} = (1 + r_{period})^c - 1
	% \label{eq:}
\end{equation}

One major limitation of annualizing returns is the implicit assumption that earned returns can be repeated.
% subsection Compounding: Daily, Monthly, Quarterly, and Annual Returns (end)

\subsection{Continuous Compounding} % (fold)
% \label{sub:Continuous Compounding}
The more frequently we compound returns, the higher the annual return will be; however, there is a limit.
\begin{itemize}
	\item \textbf{continuously compounded return}: 
	\begin{equation}
		\tilde{r}_{t, t+1} = ln \frac{P_{t+1}}{P_t} = ln(1 + r_{t, t+1})
		% \label{eq:}
	\end{equation}
\end{itemize}
The additive property of continuously compounded returns provides calculation advantages:
\begin{align*}
	\frac{P_T}{P_0} & = \frac{P_T}{P_{T-1}} \cdot \frac{P_{T-1}}{P_{T-2}} ... \frac{P_2}{P_1} \cdot \frac{P_1}{P_0} \\
	\Rightarrow ln(\frac{P_T}{P_0}) &= ln(\frac{P_T}{P_{T-1}}) + ln(\frac{P_{T-1}}{P_{T-2}}) + ... + ln(\frac{P_2}{P_1}) + ln(\frac{P_1}{P_0}) \\
	\Rightarrow \tilde{r}_{0, T} &=  \tilde{r}_{T-1, T} + \tilde{r}_{T-2, T-1} + ... + \tilde{r}_{2, 1} + \tilde{r}_{1, 0} = \sum_{t=1}^{T} \tilde{r}_{t-1, t}
\end{align*}
% subsection Continuous Compounding (end)

\section{COMMON RETURN MEASURES}
\begin{center}
\begin{tabular}{p{0.06\linewidth} | p{0.88\linewidth}}
	b & describe, compare, and interpret required rates of return, risk-free rates, risk premia, and inflation \\
\end{tabular}
\end{center}

\subsection{The Required Rate of Return} % (fold)
% \label{sub:The Required Rate of Return}
\begin{itemize}
	\item \textbf{opportunity cost}: that incurs when delaying consumption, or the value that investors forgo by choosing a course of action.
	\item \textbf{required rate of return}: the minimum rate of return an investor must receive to accept the investment.
\end{itemize}
Components of Required Rate of Return:
\begin{itemize}
	\item \textbf{risk-free rate}: is the return on an investment without any default or reinvestment risk.
	\begin{itemize}
		\item[] It refers to returns on short-term government bonds, such as the 30-day US Treasury bill. The bonds have a short maturity, which reduces their sensitivity to interest rates changes.
		\item[] \textbf{real risk-free rate} reflects the time preferences of individuals for current versus future real consumption.
		\begin{equation}
				r_{rf} = \frac{1+r_f}{1+infl} - 1
				\text{ or quick approximation: } r_{rf} = r_f - infl
			% \label{eq:}
		\end{equation}
	\end{itemize}
	\item Total returns for any asset can then be approximated as:
	\begin{equation}
	r = r_{rf} + inf + \sum_{j=1}^{n}f_j \times rp_j	
		% \label{eq:}
	\end{equation}
	\begin{itemize}
		\item[] $n$ is the number of risk premia
		\item[] $f_{ij}$ is the exposure of security $i$ to the $j^{th}$ risk factor
		\item[] $rp_j$ is the specific risk premium associated with the $j^{th}$ risk factor
		\item[] $\sum_{j=1}^{n}f_j \times rp_j$ is the sum of all risk premia.
	\end{itemize}
\end{itemize}
Examples of risk premia for bonds:
\begin{itemize}
	\item \textbf{Default risk premium}: the additional return for buying a riskier corporate bond instead of a less-risky government bond, which depends on the creditworthiness of the issuer.
	\item \textbf{Liquidity premium}: the additional return for holding a bond that may be harder to sell quickly or without impacting the price.
	\item \textbf{Maturity risk premium}: the additional return for holding more interest rate sensitive, longer maturity bonds.
\end{itemize}
Examples of risk premia for stocks:
\begin{itemize}
	\item \textbf{Size risk premium}: the return spread between small and large market-capitalization stocks.
	\item \textbf{Value risk premium}: the return spread between companies with high book-to-market and low book-to-market ratios.
\end{itemize}
% subsection The Required Rate of Return (end)
\subsection{Excess Returns} % (fold)
% \label{sub:Excess Returns}
\textbf{Excess returns}: the difference between an asset's return and the return on some reference rate or benchmark.
\begin{equation}
		\dot{r}_i = \frac{1+r_i}{1+r_B} - 1
		\text{ or quick approximation: } \dot{r}_i = r_i - r_B
	% \label{eq:}
\end{equation}
% subsection Excess Returns (end)

\subsection{Real returns} % (fold)
% \label{sub:Real returns}
\textbf{Real returns} are inflation-adjusted returns.
\begin{equation}
		r_{real} = \frac{1+r_i}{1+infl} - 1
	% \label{eq:}
\end{equation}
% subsection Real returns (end)

\subsection{Gross and Net Return} % (fold)
% \label{sub:Gross and Net Return}
\textbf{Gross return} = Total return - Transaction costs
\begin{itemize}
	\item[] an appropriate measure for evaluating and comparing the investment skill of managers,
	\item[] excludes indirect expenses related to the management and administration of an investment
\end{itemize}
\textbf{Net return} = Gross return - Management fees - Administrative expenses 
\begin{itemize}
	\item[] the return the investor actually earns
\end{itemize}
% subsection Gross and Net Return (end)

\subsection{Pre-Tax and After-Tax Nominal Return} % (fold)
% \label{sub:Pre-Tax and After-Tax Nominal Return}
\textbf{After-tax nominal return}
\begin{equation}
	r_{net-tax} = r_{price} \times (1-tax_{capital gains}) + r_{capital} \times (1-tax_{capital distribution})
	% \label{eq:}
\end{equation}
Capital gains tax is a tax on a positive price return, which can be offset for tax purposes against capital losses. Two forms of capital gains tax rate:
\begin{itemize}
	\item[] Short-term capital gains are levied on capital gains on short-term holdings
	\item[] Long-term capital gains are levied on longer-term holdings and receive preferential tax treatment in a number of countries
\end{itemize}
Forms of capital distribution tax rate:
\begin{itemize}
	\item[] Interest income is taxed as ordinary income in most countries. 
	\item[] Dividend income may be taxed as ordinary income, may have a lower tax rate, or may be exempt from taxes depending on the country and the type of investor.
\end{itemize}
Factors affect tax components:
\begin{itemize}
	\item[] The investor's specific income sources and overall tax rates (marginal tax rate)
	\item[] The time period the investor holds the investment (long-term versus short-term)
	\item[] The type of account in which the investor holds the asset (tax-exempt, tax-deferred, or taxable)
	\item[] The different tax systems under which the investor is liable to pay taxes (local, municipal, county, state, regional, and national or federal taxes)
\end{itemize}
% subsection Pre-Tax and After-Tax Nominal Return (end)

\subsection{Leveraged Returns} % (fold)
% \label{sub:Leveraged Returns}
Using leverage to increase the total amount of investible assets is a common feature of investments, whether through borrowing or derivatives.
Leverage increases investible capital, and the larger investment magnifies returns, and consequently, risks.
\begin{equation}
	R_L = \frac{Net Return}{Equity} = \frac{R_A \times A - R_D \times D}{E} = R_A + \frac{D}{E}(R_A-R_D)
	% \label{eq:}
\end{equation}
As long as the $r_p > r_d$, then $r_p$ will be positive and the leverage adds value to the portfolio.
% subsection Leveraged Returns (end)

